\chapter*{Preface}

\addcontentsline{toc}{chapter}{Preface}

Topology is the study of geometric properties and spatial relations that remain invariant under continuous deformations. For anyone stepping into modern mathematics, topology serves as a vital bridge connecting classical analysis, abstract algebra and geometry. These notes are born out of a desire to systematically organize, refine and document my personal journey through this profound and beautiful landscape.

The material presented in this book is structured into three main phases, reflecting a progression from intuitive frameworks to abstract structures and ultimately to algebraic invariants:

\begin{itemize}
    \item \textbf{General and Point-Set Topology (Chapters 1--3, 7):} We begin on the familiar ground of metric spaces, establishing the foundational intuition of open sets and continuity, before lifting these concepts into general topological spaces. Core topological properties---such as compactness, connectivity, separation axioms, and regional properties---are explored thoroughly.
    \item \textbf{Categorical Language and Geometric Constructions (Chapters 4--6, 9, 14):} Uniquely, these notes introduce the preliminaries of Category Theory at an early stage. Reframing products, disjoint unions, pullbacks, and limits through universal properties provides a unifying perspective. This categorical and combinatorial machinery is then applied to study quotient spaces, group actions, simplicial complexes, topological graphs, and the foundational definitions of topological manifolds.
    \item \textbf{An Invitation to Algebraic Topology (Chapters 8, 10--13):} The final arc transitions into algebraic topology, where spaces are investigated via algebraic invariants. Starting with the concept of homotopy and simplicial approximation, we develop the machinery of the fundamental group ($\Pi_1$). We explore practical computations using the Seifert--van Kampen Theorem and conclude with an elegant application to Knot Theory and knot groups.
\end{itemize}

These notes represent a compilation of insights gathered from standard literature, lecture courses, and personal reflections. Mathematics is an ongoing, collaborative pursuit; despite meticulous proofreading, errors, typos, and expository gaps are bound to exist. I warmly welcome any feedback, corrections, or mathematical discussions from fellow readers.


\begin{flushright}
    \textit{Wang Yuyao}
\end{flushright}